\section*{Contoh}
\addcontentsline{toc}{section}{Contoh}

\begin{contoh}[Aturan ekspresi Bison]
Aturan
\begin{lstlisting}[language=Bison]
expr : expr T_PLUS expr
     | expr T_STAR expr
     | T_INT_LIT
     ;
\end{lstlisting}
dapat menimbulkan konflik shift-reduce karena precedence tidak
dideklarasikan. Dengan \code{\%left T\_PLUS} dan
\code{\%left T\_STAR}, Bison menyelesaikan konflik sesuai prioritas
perkalian.
\end{contoh}

\begin{contoh}[Membangun simpul AST]
Aksi
\begin{lstlisting}[language=Bison]
expr : expr T_PLUS expr
     {
         $$ = ast_new(AST_BINARY, $2.token);
         $$->left = $1;
         $$->right = $3;
     }
     ;
\end{lstlisting}
membuat simpul \code{BINARY} dengan token operator dari \code{\$2},
anak kiri dari \code{\$1}, dan anak kanan dari \code{\$3}.
\end{contoh}

\begin{contoh}[AST untuk \(a + b * c\)]
Aksi \code{term} membangun simpul \code{*} untuk \(b * c\). Aksi
\code{additive} membangun simpul \code{+} dengan anak \code{a} dan
simpul \code{*}. Hasilnya adalah AST pada
\figref{fig:ch13-ast}.
\end{contoh}

\begin{contoh}[Precedence dan associativity]
Deklarasi
\begin{lstlisting}[language=Bison]
%left T_PLUS T_MINUS
%left T_STAR T_SLASH T_PERCENT
\end{lstlisting}
membuat \code{*} lebih tinggi daripada \code{+}. Keduanya
kiri-asosiatif, sehingga \(a - b - c\) diurai sebagai
\((a - b) - c\). Untuk operator kanan-asosiatif, gunakan
\code{\%right}.
\end{contoh}

\begin{contoh}[Location tracking]
Dengan \code{\%locations}, Bison menyediakan \code{@1}, \code{@2},
dan \code{@\$} untuk lokasi simbol. Aksi dapat menyalin
\code{@\$.first\_line} ke simpul AST, sehingga diagnostik semantik
dapat menyebut baris yang tepat.
\end{contoh}
